\begin{afterword}
\summaryhead

The post asks whether the proverb ``Outside the laboratory, scientists are no wiser than anyone else'' is true. The author thinks it is partly true, and finds that alarming. A shepherd trained to count sheep who cannot count apples has not understood counting. An economist who buys a lottery ticket every week may not understand expected utility. Richard Feynman's students in Brazil had memorized optics and did not recognize it in light reflected from water.

Experiments matter because to map a territory you must interact with it causally. A scientist who believes in a spirit world and says observation cannot settle the question has learned the ritual without its reason. If spirits speak to the heart, that is a causal interaction and should be weighed like any evidence. A scientist who accepts that spiritual matters lie beyond observation treats experiment as a social convention. Newton showed that the planets obey the same laws as falling apples, and reality is one process under simple laws. The author hopes that rationalists can apply the rules in everyday life as well as in the laboratory.

\responsehead

The post's general argument is clear and holds together. Evidence works the same way everywhere, and a report of inner experience is itself evidence, to be weighed against other causes and against a prior, not set aside. The problem is the claims about scientists that the post draws from it.

The premise, that the proverb ``does appear true to some degree,'' comes with no evidence. On the post's own example, religious belief, the data that exist fit both the premise and the post's hedge: a 2009 survey found American scientists about half as likely as the public to believe in God or a higher power, and about half of them still did.

The diagnosis of such scientists is based on one imagined case. From a hypothetical reply, the post concludes that the scientist ``literally doesn't know why you have to look at things.'' It then offers, as its own view (``to me''), the claims that such scientists understand experiment only as a social convention; that they ``probably never did understand'' why the rules work and ``don't like to be constrained by evidence''; and that this makes the author ``wonder'' whether to trust them even in their own field. Each step is hedged, and none is backed by anything about how real scientists who hold such beliefs reason.

A smaller example has a similar gap. The economist who buys a weekly lottery ticket is offered as someone who may not understand expected utility. Whether the ticket fails that test depends on how much the anticipation is worth, the point in dispute in ``Lotteries: A Waste of Hope.''

In short: a clear statement that the rules of evidence are the same inside and outside the laboratory, together with a diagnosis of religious scientists based on an imagined case and on no evidence about real ones.
\end{afterword}
